\documentclass[11pt,a4paper]{article}
\usepackage[margin=2.5cm, headheight=15pt]{geometry}
\usepackage{amsmath,amssymb,amsthm}
\usepackage{mathtools}
\usepackage{microtype}
\usepackage{hyperref}
\usepackage{cleveref}
\usepackage{enumitem}
\usepackage{mdframed}
\usepackage{xcolor}
\usepackage{titlesec}
\usepackage{booktabs}
\usepackage{tabularx}
\usepackage{fancyhdr}

\definecolor{navyblue}{RGB}{26,58,92}
\definecolor{midblue}{RGB}{44,95,138}
\definecolor{lightblue}{RGB}{232,240,248}
\definecolor{verylightblue}{RGB}{248,251,254}
\definecolor{proofgray}{RGB}{242,242,242}
\definecolor{warncolor}{RGB}{255,245,230}
\definecolor{warnborder}{RGB}{200,100,0}
\definecolor{opencolor}{RGB}{235,245,235}
\definecolor{openborder}{RGB}{0,120,0}

\hypersetup{colorlinks=true,linkcolor=midblue,citecolor=midblue,urlcolor=midblue,
	pdftitle={SSP Part IV: GTSP-MTZ Formulation},bookmarksnumbered=true}

\pagestyle{fancy}\fancyhf{}
\fancyhead[L]{\small\color{navyblue}\textit{Part IV: GTSP-MTZ Formulation for the Uniform SSP}}
\fancyhead[R]{\small\color{navyblue}\thepage}
\renewcommand{\headrulewidth}{0.4pt}

\titleformat{\section}{\large\bfseries\color{navyblue}}{}{0em}{}[\color{navyblue}\titlerule]
\titleformat{\subsection}{\normalsize\bfseries\color{midblue}}{}{0em}{}
\titleformat{\subsubsection}{\normalsize\bfseries\itshape\color{midblue}}{}{0em}{}
\titlespacing{\section}{0pt}{18pt}{8pt}
\titlespacing{\subsection}{0pt}{12pt}{5pt}

\newmdenv[backgroundcolor=verylightblue,linecolor=navyblue,linewidth=1.2pt,
leftline=true,rightline=false,topline=false,bottomline=false,
innerleftmargin=12pt,innerrightmargin=8pt,innertopmargin=8pt,innerbottommargin=8pt,
skipabove=10pt,skipbelow=6pt]{thmbox}

\newmdenv[backgroundcolor=proofgray,linecolor=midblue,linewidth=0.8pt,
leftline=true,rightline=false,topline=false,bottomline=false,
innerleftmargin=10pt,innerrightmargin=8pt,innertopmargin=6pt,innerbottommargin=6pt,
skipabove=6pt,skipbelow=6pt]{proofbox}

\newmdenv[backgroundcolor=lightblue,linecolor=navyblue,linewidth=1.5pt,roundcorner=3pt,
innerleftmargin=14pt,innerrightmargin=14pt,innertopmargin=10pt,innerbottommargin=10pt,
skipabove=12pt,skipbelow=8pt]{formbox}

\newmdenv[backgroundcolor=warncolor,linecolor=warnborder,linewidth=1.2pt,
leftline=true,rightline=false,topline=false,bottomline=false,
innerleftmargin=12pt,innerrightmargin=8pt,innertopmargin=8pt,innerbottommargin=8pt,
skipabove=10pt,skipbelow=6pt]{warnbox}

\newmdenv[backgroundcolor=opencolor,linecolor=openborder,linewidth=1.2pt,
leftline=true,rightline=false,topline=false,bottomline=false,
innerleftmargin=12pt,innerrightmargin=8pt,innertopmargin=8pt,innerbottommargin=8pt,
skipabove=10pt,skipbelow=6pt]{openproblem}

\theoremstyle{plain}
\newtheorem{theorem}{Theorem}[section]
\newtheorem{lemma}[theorem]{Lemma}
\newtheorem{proposition}[theorem]{Proposition}
\theoremstyle{definition}
\newtheorem{definition}[theorem]{Definition}
\newtheorem{remark}[theorem]{Remark}

\DeclareMathOperator{\OPT}{OPT}
\newcommand{\OMTZ}{\OPT_{\mathrm{MTZ}}}
\newcommand{\OSSP}{\OPT_{\mathrm{SSP}}}

\setlist[itemize]{noitemsep,topsep=4pt}
\setlist[enumerate]{noitemsep,topsep=4pt}

%-------------------------------------------------------------------------------
\begin{document}
%-------------------------------------------------------------------------------

\begin{center}
	{\color{navyblue}\rule{\textwidth}{2pt}}\\[0.4cm]
	{\LARGE\bfseries\color{navyblue}The Job Sequencing and Tool Switching Problem}\\[0.2cm]
	{\large\bfseries\color{navyblue}Part IV: A Cluster-Level MTZ Formulation for the Non-Compact SSP}\\[0.1cm]
	{\large\itshape\color{midblue}Why it is not exact (job domination), the safe formulations to use, and the LP-strength trade-off}\\[0.3cm]
	{\color{navyblue}\rule{\textwidth}{2pt}}
\end{center}

\vspace{0.3cm}

\begin{mdframed}[backgroundcolor=lightblue,linecolor=navyblue,linewidth=1pt,
  innerleftmargin=14pt,innerrightmargin=14pt,innertopmargin=8pt,innerbottommargin=8pt]
\small\textbf{Scope and status.}
This note proposes replacing the exponentially many subtour-elimination constraints of
the non-compact SSP formulation (Part~II) by a polynomial ($O(n^2)$) family of
cluster-level Miller--Tucker--Zemlin (MTZ) constraints. Two facts must be stated up front,
because they determine how the formulation may legitimately be used.
\begin{enumerate}[label=\textbf{(\arabic*)}]
\item \textbf{The model is \emph{not} exact (now settled).} The aggregate MTZ constraint does
not merely forbid subtours: it forbids any solution that traverses an ordered cluster pair
twice and, through the job-indexed potentials $u_j$, any solution whose induced precedence
digraph has a cycle (\cref{sec:exchange}). It thus optimises over a strict subset of
schedules, so $\OMTZ\ge\OSSP$ --- a heuristic upper bound. We previously left exactness
(the ``exchange argument'') open; it is now \textbf{resolved negatively}: there is a $b=3$
counterexample (\cref{thm:notexact}) on which $\OMTZ=2>1=\OSSP$, and the failure is generic
(it is triggered by \emph{job domination} $T_i\subseteq T_j$, which is ubiquitous).
\textbf{Therefore this formulation must not be used as an exact model for column generation.}
The safe exact alternatives --- per-configuration MTZ and lazy SEC/PCTSP cuts --- are given
in \cref{sec:safe}.
\item \textbf{The motivation is a trade-off, not a free win.} Removing row generation costs
LP strength: MTZ-type subtour constraints give a relaxation no tighter than SEC, and
aggregation to clusters weakens it further (\cref{rem:lpweak}). Whether the saved separation
outweighs the looser bound inside branch-and-price is an \emph{open computational question}
(\cref{sec:complexity}); this note claims no net improvement.
\end{enumerate}
\end{mdframed}

\vspace{0.3cm}
\tableofcontents
\newpage

%-------------------------------------------------------------------------------
\section{Motivation: The Double-Generation Problem and the MTZ Trade-off}
\label{sec:motivation}
%-------------------------------------------------------------------------------

The non-compact configuration ILP of Part~II
(\texttt{02\_noncompact\_formulation.tex}) has two simultaneously exponential components:
\begin{itemize}
\item \textbf{Exponential columns}: $|V|=\binom{|T|}{b}$ configuration/arc variables,
      handled by \textbf{column generation} (a set-union knapsack pricing problem);
\item \textbf{Exponential rows}: subtour-elimination constraints (SEC/GSEC), handled by
      \textbf{row generation} (separation).
\end{itemize}
Running both at once (``double generation'') is awkward: a newly priced column can violate
previously satisfied subtour cuts, forcing repeated separation passes and complicating the
branch-and-price bookkeeping \cite{Barnhart1998,Lubbecke2005}.

\medskip\noindent\textbf{The proposal.} Replace the SECs by $O(n^2)$ cluster-level MTZ
constraints \cite{Miller1960}. Row generation then disappears; only column generation
remains. The constraint-side comparison is:
\begin{center}
	\renewcommand{\arraystretch}{1.3}
	\begin{tabular}{lcc}
		\toprule
		& \textbf{SEC-based (Part II)} & \textbf{MTZ-based (this note)} \\
		\midrule
		Variables ($x_{ij}$) & $|V|^2$ (exponential) & $|V|^2$ (exponential) \\
		Subtour constraints & $\sim 2^{|V|}$ & $n(n-1)$ (polynomial) \\
		Coverage constraints & $n$ & $2n$ \\
		Column generation? & Yes & Yes \\
		Row generation? & \textbf{Yes} & \textbf{No} \\
		LP relaxation strength & \textbf{stronger} & weaker (\cref{rem:lpweak}) \\
		Captures every SSP optimum? & Yes & \emph{Conditional} (\cref{thm:correct}) \\
		\bottomrule
	\end{tabular}
\end{center}

\begin{warnbox}
\textbf{What is gained and what is paid.}
The last two rows are the crux and were glossed over in earlier drafts. Removing row
generation is genuine. But (i)~MTZ subtour constraints give a weaker LP relaxation than SEC,
and the cluster aggregation weakens it further (\cref{rem:lpweak}); and (ii)~the aggregate
constraint additionally \emph{restricts} the integer feasible set, so the model can miss SSP
optima unless the exchange argument holds (\cref{sec:exchange}). The proposal thus trades a
tight, exact SEC model (much separation) for a cheap, conditionally exact MTZ model (no
separation, weaker bound). Which wins inside branch-and-price is empirical; we flag it in
\cref{sec:complexity}. The Desrochers--Laporte lifting (\cref{rem:dl}) recovers part of the
lost LP strength at the same constraint count.
\end{warnbox}

%-------------------------------------------------------------------------------
\section{Notation}
\label{sec:notation}
%-------------------------------------------------------------------------------

\begin{itemize}
	\item $J=\{1,\dots,n\}$: jobs;\quad $T$: tools;\quad $b$: magazine capacity.
	\item $T_j\subseteq T$, $|T_j|\le b$: mandatory tools of job $j$.
	\item $V$: all valid configurations --- subsets $\mathcal{C}\subseteq T$ with
	$|\mathcal{C}|=b$. $|V|=\binom{|T|}{b}$, exponential in general.
	\item Configuration cluster of job $j^{*}$:
	$\displaystyle H_{j^{*}}=\bigl\{\mathcal{C}\in V:\ T_{j^{*}}\subseteq\mathcal{C}\bigr\}$,
	the configurations that can serve $j^{*}$. Clusters may overlap.
	\item Switching cost between $i,j\in V$:\ $d_{ij}=b-|i\cap j|=|i\setminus j|$.
	\item Dummy source/sink $s,t\notin V$ for the path formulation (zero-cost arcs to/from
	all of $V$).
\end{itemize}

%-------------------------------------------------------------------------------
\section{The Cluster-Level GTSP-MTZ Formulation}
\label{sec:formulation}
%-------------------------------------------------------------------------------

Let $V^{+}=V\cup\{s,t\}$ and $x_{ij}\in\{0,1\}$ for $(i,j)\in V^{+}\times V^{+}$,
$i\ne j$, with position variables $u_{j'}$ for each job $j'\in J$.

\begin{formbox}
	\textbf{GTSP-MTZ formulation for the SSP (path version)}

	\smallskip
	\textbf{Objective:}
	\begin{equation}\label{eq:obj}
		\min \quad \sum_{i\in V}\sum_{j\in V,\,j\neq i} d_{ij}\,x_{ij}
	\end{equation}
	\textbf{Path start/end:}
	\begin{align}
		\sum_{j\in V} x_{sj}&=1, \label{eq:source}\\
		\sum_{i\in V} x_{it}&=1. \label{eq:sink}
	\end{align}
	\textbf{Degree (each config used at most once):}
	\begin{align}
		\sum_{j\in V^{+}\setminus\{i\}} x_{ij}&\le 1 \quad \forall i\in V,\label{eq:outdeg}\\
		\sum_{i\in V^{+}\setminus\{j\}} x_{ij}&\le 1 \quad \forall j\in V.\label{eq:indeg}
	\end{align}
	\textbf{Flow conservation:}
	\begin{equation}\label{eq:flow}
		\sum_{j\in V^{+}\setminus\{i\}} x_{ij}=\sum_{j\in V^{+}\setminus\{i\}} x_{ji}
		\quad\forall i\in V.
	\end{equation}
	\textbf{Cluster coverage --- exit and entry of each job cluster:}
	\begin{align}
		\sum_{i\in H_{j^{*}}}\sum_{j\in V\setminus H_{j^{*}}} x_{ij}&\ge 1
		\quad\forall j^{*}\in J,\label{eq:exit}\\
		\sum_{i\in V\setminus H_{j^{*}}}\sum_{j\in H_{j^{*}}} x_{ij}&\ge 1
		\quad\forall j^{*}\in J.\label{eq:entry}
	\end{align}
	\textbf{Cluster-level MTZ subtour elimination:}
	\begin{equation}\label{eq:MTZ}
		u_{j'}-u_{j''}+n\!\!\sum_{i\in H_{j'}}\sum_{j\in H_{j''}}\! x_{ij}\ \le\ n-1
		\qquad\forall j',j''\in J,\ j'\neq j''.
	\end{equation}
	\textbf{Bounds and integrality:}
	\begin{equation}\label{eq:ubounds}
		1\le u_{j'}\le n\ \ \forall j'\in J,\qquad x_{ij}\in\{0,1\}\ \ \forall i,j\in V^{+}.
	\end{equation}
\end{formbox}

%-------------------------------------------------------------------------------
\section{Mathematical Verification}
\label{sec:verify}
%-------------------------------------------------------------------------------

\subsection{Objective and coverage}

The objective \eqref{eq:obj} sums $d_{ij}=b-|i\cap j|$ over selected arcs; since
$|i|=|j|=b$, $|i\setminus j|=b-|i\cap j|$ equals both the removals and (by the uniform
model) the insertions, so \eqref{eq:obj} counts tool switches.

\begin{thmbox}
	\begin{proposition}[Coverage $\Leftrightarrow$ job feasibility]
		\label{prop:coverage}
		In any feasible solution, every job $j^{*}$ is served by some visited configuration
		$\mathcal{C}\in H_{j^{*}}$.
	\end{proposition}
\end{thmbox}
\begin{proofbox}
	\begin{proof}
		The selected arcs form an $s$--$t$ path through configurations. Constraint
		\eqref{eq:exit} forces an arc leaving $H_{j^{*}}$ and \eqref{eq:entry} an arc entering
		it; on a connected path the existence of both implies the path passes through at least
		one node of $H_{j^{*}}$, which serves $j^{*}$.
	\end{proof}
\end{proofbox}

\subsection{The exchange argument: the central open problem}
\label{sec:exchange}

The aggregate MTZ constraint \eqref{eq:MTZ} carries a hidden assumption, which is the
heart of this note.

\begin{warnbox}
	\textbf{Why \eqref{eq:MTZ} \emph{restricts} rather than relaxes.}
	The sum $\sum_{i\in H_{j'}}\sum_{j\in H_{j''}} x_{ij}$ counts \emph{all} arcs from cluster
	$H_{j'}$ to cluster $H_{j''}$. If this equals $k\ge 2$, \eqref{eq:MTZ} forces
	$u_{j''}\ge u_{j'}+n(k-1)+1$, which with $1\le u\le n$ is infeasible for $k\ge2$. Hence the
	formulation \emph{forbids} every solution that traverses some ordered cluster pair more than
	once. Because it removes solutions (rather than adding fractional ones), it is a
	\emph{restriction}: it can only \emph{raise} the optimum, never lower it. It is exact only
	if an optimal SSP solution never needs such a repeated cluster pair.
\end{warnbox}

\begin{thmbox}
	\begin{definition}[Cluster-simple path]
		\label{def:cluster-simple}
		A configuration path $C_{(1)},\dots,C_{(m)}$ is \emph{cluster-simple} if for every
		ordered pair $(j',j'')$, $j'\ne j''$,
		$\sum_{l=1}^{m-1}\mathbf{1}[\,C_{(l)}\in H_{j'},\,C_{(l+1)}\in H_{j''}\,]\le 1$.
	\end{definition}
\end{thmbox}

\begin{thmbox}
	\begin{proposition}[Exact representability condition]
	\label{prop:representable}
	A configuration path is representable by \eqref{eq:MTZ} (for some $u\in[1,n]$) \emph{iff}
	(i)~every ordered cluster pair is traversed at most once (\cref{def:cluster-simple}), and
	(ii)~the precedence digraph $D$ --- edge $j'\!\to\!j''$ whenever some arc runs from a config
	in $H_{j'}$ to one in $H_{j''}$ --- is \emph{acyclic}.
	\end{proposition}
\end{thmbox}
\begin{proofbox}
	\begin{proof}
	Two arcs on a pair give sum $\ge2$ and infeasibility (above), forcing~(i). Given~(i),
	\eqref{eq:MTZ} reduces to $u_{j'}<u_{j''}$ per edge of $D$, solvable in $[1,n]$ iff $D$ is
	acyclic (topological order). \qed
	\end{proof}
\end{proofbox}

\begin{thmbox}
	\begin{theorem}[The cluster-MTZ formulation is \emph{not} exact]
	\label{thm:notexact}
	There is a $b=3$ instance with $\OMTZ>\OSSP$; the cluster-MTZ model excludes some optimal
	SSP solutions and is in general a strict over-estimate of $\OSSP$.
	\end{theorem}
\end{thmbox}
\begin{proofbox}
	\begin{proof}
	Let $b=3$, $T_1=\{1,4\}$, $T_2=\{1,3,4\}$, $T_3=\{1\}$, $T_4=\{1,2,4\}$
	($U=\{1,2,3,4\}$). The only config covering $T_2$ is $\{1,3,4\}$ and the only one covering
	$T_4$ is $\{1,2,4\}$, so every covering path uses both; they already cover all four jobs and
	$d(\{1,3,4\},\{1,2,4\})=1$, hence $\OSSP=1$ (one arc). On that arc both $T_3=\{1\}$ and
	$T_1=\{1,4\}$ are covered at \emph{both} endpoints (as $\{1\}\subseteq\{1,4\}$ lie in both
	configs), so $D$ contains $1\!\to\!3$ and $3\!\to\!1$: a $2$-cycle. By
	\cref{prop:representable} no cost-$1$ path is representable; the cheapest representable
	covering path costs $2$, so $\OMTZ=2>1$. (Brute-force verified; $19$ counterexamples in
	$\approx7.9$k instances, \texttt{plans-genai/\_verification/}.) \qed
	\end{proof}
\end{proofbox}

\begin{warnbox}
	\textbf{The failure is generic: job domination.} Whenever $T_i\subseteq T_j$, every config
	covering $j$ also covers $i$, so any arc with such a config on both ends puts $\{i,j\}$ on
	both sides and creates the $2$-cycle $i\!\leftrightarrow\!j$ in $D$. Domination underlies the
	classical SSP dominance rules and is ubiquitous, so cluster-MTZ fails on a large fraction of
	realistic instances. \textbf{It is not a safe basis for an exact column-generation solver};
	use the alternatives of \cref{sec:safe}. The propositions below give the (restrictive)
	classes on which cluster-MTZ \emph{does} stay exact.
\end{warnbox}

\begin{thmbox}
	\begin{proposition}[Holds for pairwise-incompatible instances]
		\label{prop:exchange_disjoint}
		If every pair of distinct jobs $j',j''$ has $|T_{j'}\cup T_{j''}|>b$ (so
		$H_{j'}\cap H_{j''}=\emptyset$), then every feasible configuration path --- in
		particular every optimal one --- is cluster-simple.
	\end{proposition}
\end{thmbox}
\begin{proofbox}
	\begin{proof}
		With $H_{j'}\cap H_{j''}=\emptyset$, each configuration lies in at most one cluster, so
		each arc maps to a unique ordered cluster pair. The degree constraints
		\eqref{eq:outdeg}--\eqref{eq:indeg} let each configuration appear once, so no ordered
		cluster pair receives two arcs.
	\end{proof}
\end{proofbox}

\begin{thmbox}
	\begin{proposition}[Holds for the $6$-job ring]
		\label{prop:exchange_ring}
		The $6$-ring of Part~V (\texttt{05\_jgp\_ssp\_gap\_analysis.tex}, Example~3.1) has a
		cluster-simple optimal solution.
	\end{proposition}
\end{thmbox}
\begin{proofbox}
	\begin{proof}
		With $T_{j_k}=\{t_k,t_{k+1\bmod 6}\}$, $b=3$, the optimal sequence
		$j_1\to\cdots\to j_6$ ($\OSSP=3$, certified in Part~V) uses the four configurations
		$\mathcal{C}_1=\{t_1,t_2,t_3\}$, $\mathcal{C}_2=\{t_1,t_3,t_4\}$,
		$\mathcal{C}_3=\{t_1,t_4,t_5\}$, $\mathcal{C}_4=\{t_1,t_5,t_6\}$, lying in clusters
		$\mathcal{C}_1\in H_{j_1}\cap H_{j_2}$, $\mathcal{C}_2\in H_{j_3}$,
		$\mathcal{C}_3\in H_{j_4}$, $\mathcal{C}_4\in H_{j_5}\cap H_{j_6}$. The three arcs hit
		exactly the ordered cluster pairs $(j_1,j_3),(j_2,j_3),(j_3,j_4),(j_4,j_5),(j_4,j_6)$,
		each once; all other $25$ ordered pairs are unused. With
		$u_{j_1}=u_{j_2}=1,u_{j_3}=2,u_{j_4}=3,u_{j_5}=u_{j_6}=4$ one checks
		$u_{j'}-u_{j''}+6\cdot(\text{pair-sum})\le 5$ for all $30$ ordered pairs.
	\end{proof}
\end{proofbox}

\begin{remark}[Why these classes are safe but the general case is not]
	Both classes avoid the domination $2$-cycle of \cref{thm:notexact}. In the disjoint-tool
	case each configuration lies in a single cluster, so an arc induces exactly one ordered
	cluster pair and $D$ is a simple path (acyclic). In the ring each optimal configuration
	serves at most the two ring-adjacent jobs and the optimal sequence visits clusters
	contiguously, so $D$ is again acyclic. The general case fails precisely when a configuration
	covers a dominated pair $T_i\subseteq T_j$ on both ends of an arc (\cref{thm:notexact}); no
	hypothesis rules that out for arbitrary instances.
\end{remark}

\subsection{MTZ eliminates cluster-level subtours}

\begin{thmbox}
	\begin{proposition}[No cluster-level subtours]
		\label{prop:MTZ}
		If $\sum_{i\in H_{j'}}\sum_{j\in H_{j''}} x_{ij}\le 1$ for all $(j',j'')$, then no
		directed cycle visits configurations drawn only from clusters $\{H_{j'}:j'\in S\}$ for
		any $S\subsetneq J$.
	\end{proposition}
\end{thmbox}
\begin{proofbox}
	\begin{proof}
		Suppose such a cycle exists. Summing \eqref{eq:MTZ} over its consecutive cluster pairs,
		the telescoping $u$-differences cancel and the arc term contributes $n|S|$ (one per
		pair), giving $n|S|\le|S|(n-1)$, i.e.\ $n\le n-1$, a contradiction.
	\end{proof}
\end{proofbox}

\begin{warnbox}
	\textbf{Intra-cluster subtours.} \eqref{eq:MTZ} does not forbid a cycle among several
	configurations of the \emph{same} cluster $H_{j^{*}}$. Such a cycle serves $j^{*}$
	redundantly at positive cost, so it never occurs in an integer optimum (it may occur in
	the LP relaxation). A full guarantee needs per-configuration MTZ
	$v_i-v_j+|V|x_{ij}\le|V|-1$, which reintroduces $O(|V|^{2})$ (exponential) constraints
	and is therefore not used.
\end{warnbox}

\subsection{Path structure and conditional correctness}

\begin{thmbox}
	\begin{proposition}[Path structure]
		\label{prop:path}
		Constraints \eqref{eq:source}--\eqref{eq:flow} and the degree bounds make the selected
		arcs a single directed $s$--$t$ path through a subset of $V$, once \eqref{eq:MTZ}
		removes cluster-level cycles.
	\end{proposition}
\end{thmbox}
\begin{proofbox}
	\begin{proof}
		One arc leaves $s$, one enters $t$; flow conservation and degree $\le1$ at internal
		nodes give a disjoint union of one $s$--$t$ path and cycles. \cref{prop:MTZ} removes the
		cycles (which, lying in $V$, span clusters of $J$), leaving a single path.
	\end{proof}
\end{proofbox}

\begin{thmbox}
	\begin{theorem}[Restriction; not exact in general]
		\label{thm:correct}
		The integer model \eqref{eq:obj}--\eqref{eq:ubounds} is a \emph{restriction} of the SSP
		to \emph{representable} schedules (cluster-simple with acyclic precedence digraph,
		\cref{prop:representable}): every feasible integer solution is a feasible SSP schedule
		of equal cost, so $\OMTZ\ge\OSSP$, and $\OMTZ$ is the cheapest representable schedule.
		Equality $\OMTZ=\OSSP$ holds \emph{iff} some optimal SSP schedule is representable; this
		\textbf{fails in general} (\cref{thm:notexact}) but holds for the classes of
		\cref{prop:exchange_disjoint,prop:exchange_ring}.
	\end{theorem}
\end{thmbox}
\begin{proofbox}
	\begin{proof}
		\emph{Every model solution is a schedule (hence $\OMTZ\ge\OSSP$).} By
		\cref{prop:path,prop:coverage} an integer-feasible solution is a single $s$--$t$ path
		that serves every job; processing each job at a covering configuration along the path
		gives an SSP schedule whose switch cost equals the objective \eqref{eq:obj}. Thus a
		schedule of cost $\OMTZ$ exists, so $\OSSP\le\OMTZ$. By \cref{prop:MTZ} and the
		over-constraining identity of \cref{sec:exchange}, the feasible solutions are exactly the
		cluster-simple schedules, so $\OMTZ$ is their minimum cost.

		\emph{Conditional equality.} If some optimal SSP schedule is \emph{representable}
		(\cref{prop:representable}), collapse its repeated configurations to
		$\mathcal{C}_{(1)},\dots,\mathcal{C}_{(k)}$, set
		$x_{s\mathcal{C}_{(1)}}=x_{\mathcal{C}_{(k)}t}=1$,
		$x_{\mathcal{C}_{(l)}\mathcal{C}_{(l+1)}}=1$, and take $u$ from a topological order of the
		acyclic precedence digraph $D$. Coverage, degree and flow hold by construction;
		cluster-simplicity gives every ordered cluster-pair sum $\le1$ and acyclicity makes the
		$u$-order consistent, so \eqref{eq:MTZ} holds, giving $\OMTZ\le\OSSP$. When \emph{no}
		optimal schedule is representable (\cref{thm:notexact}), $\OMTZ>\OSSP$ strictly.
	\end{proof}
\end{proofbox}

\begin{warnbox}
	\textbf{Reading of \cref{thm:correct}.} The model always returns the cheapest
	\emph{cluster-simple} schedule, which is a genuine SSP schedule --- an \emph{upper} bound on
	$\OSSP$. When no optimum is cluster-simple, $\OMTZ>\OSSP$: the model \emph{misses} the true
	optimum. A certified \emph{lower} bound must therefore come from a genuine relaxation (the
	base degree/flow/coverage model, or the SEC formulation), \emph{not} from the
	over-constrained MTZ model. Exactness \emph{fails} in general (\cref{thm:notexact}).
\end{warnbox}

%-------------------------------------------------------------------------------
\section{Constraint Count, LP Strength, and Column Generation}
\label{sec:complexity}
%-------------------------------------------------------------------------------

\begin{center}
	\renewcommand{\arraystretch}{1.35}
	\begin{tabularx}{\textwidth}{@{}lXr@{}}
		\toprule
		\textbf{Group} & \textbf{Description} & \textbf{Count} \\
		\midrule
		\eqref{eq:source},\eqref{eq:sink} & Path start/end & 2 \\
		\eqref{eq:outdeg},\eqref{eq:indeg} & Config degree $\le1$ & $2|V|$ \\
		\eqref{eq:flow} & Flow conservation & $|V|$ \\
		\eqref{eq:exit},\eqref{eq:entry} & Cluster entry/exit coverage & $2n$ \\
		\eqref{eq:MTZ} & Cluster-level MTZ & $n(n-1)$ \\
		\eqref{eq:ubounds} & Position bounds & $2n$ \\
		\midrule
		\textbf{Total} & ($2n+n(n-1)+2n+2=(n+1)(n+2)$) & $\boldsymbol{3|V|+(n+1)(n+2)}$ \\
		\bottomrule
	\end{tabularx}
\end{center}
The $|V|$-indexed constraints are handled by column generation; the job-indexed groups
(coverage, MTZ, bounds) are $O(n^{2})$ --- polynomial, with no row generation. This is the
intended structural gain over the $O(2^{|V|})$ SECs.

\begin{remark}[LP strength: MTZ is no tighter than SEC]
	\label{rem:lpweak}
	For routing problems it is classical that compact MTZ subtour elimination yields an LP
	relaxation no tighter than --- and generally strictly weaker than --- the
	subtour-elimination (SEC/DFJ) relaxation, because MTZ enforces connectivity only through
	the potentials $u$ \cite{Applegate2006,DesrochersLaporte1991}. Two features make the
	present relaxation weaker still: the potentials are indexed by \emph{jobs} (at most $n$
	values) rather than by configurations, and the per-arc terms are \emph{aggregated} over
	$H_{j'}\times H_{j''}$, which coarsens the constraint. We therefore expect the root LP bound
	of this model to be looser than both the SEC model of Part~II and the SSPMF bound $|U|-b$ of
	Part~V; a precise polyhedral comparison for the cluster-aggregated system is left open.
\end{remark}

\begin{openproblem}
	\textbf{Net effect (open, computational).} \cref{rem:lpweak} says the root bound can only
	get looser, and \cref{thm:correct} says the integer model may miss optima. Against this, the
	model needs no separation and has $O(n^2)$ rows. Whether branch-and-price is faster overall
	--- weaker bound and a possibly larger or truncated tree, but cheaper node LPs and no cut
	management --- must be measured on the Catanzaro Tabela1C instances against the SEC/GSEC and
	SSPMF baselines. This note makes no such claim.
\end{openproblem}

\begin{remark}[Column-generation compatibility]
	\label{rem:cg}
	A newly priced arc $x_{i^{*}j^{*}}$ (with $i^{*}\in H_{j'}$, $j^{*}\in H_{j''}$) simply
	joins the existing MTZ constraint \eqref{eq:MTZ} for the pair $(j',j'')$ with coefficient
	$n$; no new constraint is created. This is the standard behaviour of column generation ---
	new variables acquire coefficients in existing rows \cite{Lubbecke2005} --- and is exactly
	what avoids the SEC double-generation issue. The pricing problem seeks the arc of most
	negative reduced cost
	$\bar d_{ij}=d_{ij}-\mu_i-\lambda_j-\sum_{j^{*}:\,i\in H_{j^{*}},\,j\notin H_{j^{*}}}\pi_{j^{*}}$.
\end{remark}

\begin{remark}[Desrochers--Laporte lifting recovers some strength]
	\label{rem:dl}
	For per-arc MTZ, Desrochers and Laporte \cite{DesrochersLaporte1991} strengthen
	$u_{j'}-u_{j''}+n\,x_{ij}\le n-1$ to
	$u_{j'}-u_{j''}+n\,x_{ij}+(n-2)x_{ji}\le n-1$ at no extra constraint count. The aggregate
	analogue replaces the arc terms by the forward and reverse cluster-pair sums. This tightens
	the relaxation toward (but not to) SEC; by how much, on SSP instances, is part of the open
	question above. A per-arc (non-aggregated) MTZ would give $O(n^{2}|V|^{2})$ constraints ---
	exponential --- and is not used.
\end{remark}

%-------------------------------------------------------------------------------
\section{Structural Elements of the Path Formulation}
\label{sec:structure}
%-------------------------------------------------------------------------------

\begin{enumerate}[label=(\arabic*)]
	\item \textbf{Dummy source/sink} \eqref{eq:source}--\eqref{eq:sink}: the SSP is a path,
	not a Hamiltonian cycle, through the configuration graph; equivalently set $d_{ts}=0$ and
	break the cycle at $(t\to s)$.
	\item \textbf{Flow conservation} \eqref{eq:flow}: every entered internal configuration is
	exited, so the path cannot stop at an interior node.
	\item \textbf{Degree bounds} \eqref{eq:outdeg}--\eqref{eq:indeg}: simple path, no
	multigraph \cite{Applegate2006}.
	\item \textbf{Cluster-level MTZ} \eqref{eq:MTZ}: $O(n^{2})$ subtour elimination in place of
	SECs, but \emph{not exact} (\cref{thm:notexact}); LP-weaker than SEC
	(\cref{rem:lpweak}) and integer-restricting. Use \cref{sec:safe} for an exact solver.
\end{enumerate}

\begin{thmbox}
	\begin{remark}[Why $n$ is the right MTZ coefficient]
		The variables $u_{j'}$ order \emph{cluster visits}; there are $\le n$ clusters, so
		$u_{j'}\in[1,n]$ and the coefficient $n$ is tight. Using $|V|$ would only weaken the
		relaxation.
	\end{remark}
\end{thmbox}

%-------------------------------------------------------------------------------
\section{What to Use Instead (for Exact Column Generation)}
\label{sec:safe}
%-------------------------------------------------------------------------------

Since the cluster-MTZ model is not exact (\cref{thm:notexact}), an exact branch-and-price
solver should use one of the following; both keep the configuration column generation intact
and only change how subtours are eliminated.
\begin{enumerate}[label=(\alph*)]
\item \textbf{Per-configuration MTZ (exact).} Give each configuration node $i\in V$ a
potential $v_i\in[1,|V|]$ and impose, for generated arcs,
\[
  v_i - v_j + |V|\,x_{ij} \le |V|-1 \qquad \forall\, i,j\in V,\ i\ne j .
\]
These are the \emph{standard} MTZ subtour constraints on the configuration graph: they
eliminate exactly the subtours and exclude no optimum --- there is no domination pathology,
because the potentials are per-configuration, not per-job, so the $2$-cycle of
\cref{thm:notexact} cannot arise. In column generation a new column $i$ adds one variable
$v_i$ and arc rows over generated nodes only: $O(|V_{\mathrm{gen}}|^{2})$ rows --- more than
the $O(n^{2})$ cluster form, but \emph{correct}. (This is the recommended choice for the
upcoming CG runs.)
\item \textbf{Lazy SEC/PCTSP cuts (exact).} Keep the Part~II degree/flow/coverage model and
separate subtour cuts on integer incumbents (the standard double-generation loop); fully
worked out in Part~II.
\end{enumerate}
\noindent The cluster-aggregated MTZ \eqref{eq:MTZ} remains usable as a fast \emph{primal
heuristic} (every feasible solution is a genuine schedule) or warm start --- just not as the
exact master. The job-indexed MTZ of the BBC master (Part~VIII) is on \emph{jobs}, not
configurations, and is the standard exact ATSP MTZ, hence unaffected.

%-------------------------------------------------------------------------------
\section{Summary}
\label{sec:summary}
%-------------------------------------------------------------------------------

\begin{itemize}
	\item \textbf{A restriction, \emph{not} exact.} The model optimises over representable
	schedules (cluster-simple with acyclic precedence digraph, \cref{prop:representable}), so
	$\OMTZ\ge\OSSP$ always but $\OMTZ>\OSSP$ is possible: there is a $b=3$ counterexample with
	$\OMTZ=2>1=\OSSP$ (\cref{thm:notexact}), and the failure is generic (job domination
	$T_i\subseteq T_j$). Exactness holds only on the restrictive classes
	\cref{prop:exchange_disjoint,prop:exchange_ring}. \textbf{For an exact CG solver, use the
	per-configuration MTZ or lazy SEC of \cref{sec:safe}}, not the cluster form.
	\item \textbf{Polynomial subtour constraints, no row generation.} $O(n^{2})$ MTZ
	constraints \cite{Miller1960} replace $\sim2^{|V|}$ SECs; new columns never trigger
	separation (\cref{rem:cg}).
	\item \textbf{At an LP-strength cost.} The MTZ relaxation is no tighter than SEC, and the
	cluster aggregation weakens it further (\cref{rem:lpweak}); Desrochers--Laporte lifting
	(\cref{rem:dl}) recovers part of it. The net branch-and-price effect is an open
	computational question (\cref{sec:complexity}).
	\item \textbf{Still column-generated.} $|V|=\binom{|T|}{b}$ remains exponential; the model
	is meant for branch-and-price \cite{Barnhart1998,Lubbecke2005}.
\end{itemize}

%-------------------------------------------------------------------------------
\section{Where This Fits: Solution Architecture}
\label{sec:fits}
%-------------------------------------------------------------------------------

\begin{center}
\renewcommand{\arraystretch}{1.3}
\begin{tabular}{lccc}
\toprule
\textbf{Formulation} & \textbf{Column gen.?} & \textbf{Row gen.?} & \textbf{Role} \\
\midrule
Compact ILPs (Part I) & No & No & exact, weak--medium LP \\
Non-compact $+$ SEC (Part II) & Yes & \textbf{Yes} & exact, strong LP \\
\textbf{MTZ (this note)} & \textbf{Yes} & \textbf{No} & restriction; heuristic UB \\
Branch-and-Benders & No & Yes (lazy) & exact, polynomial subproblem \\
\bottomrule
\end{tabular}
\end{center}

\medskip
The MTZ route removes row generation, simplifying branch-and-price
\cite{Barnhart1998,Lubbecke2005}, but inherits the classical MTZ-vs-SEC LP weakness
\cite{Applegate2006} and additionally restricts the integer set. Until the exchange argument
is settled it should be used as a heuristic that returns genuine schedules (upper bounds on
$\OSSP$), with a certified lower bound taken from a true relaxation; it must be benchmarked
against the SEC and SSPMF formulations before adoption as an exact solver.

\bibliographystyle{abbrvnat}
\bibliography{references}
\end{document}
